
\documentclass{article}
\usepackage{zref-clever}
\usepackage{../latex-course}
 \usepackage{mathtools}
 \usepackage{amssymb}

\courseweek{03}{Math Formatting -- Part 1}
\coursestudents{Sandrine Pudukkattukaran Joy (4743017) \and Felix Winter (4751930)}

\title{Inner Products and the Exponential Function}
\author{Sandrine Pudukkattukaran Joy (4743017), Felix Winter (4751930)}
\DeclarePairedDelimiter{\norm}{\lVert}{\rVert}
\DeclarePairedDelimiter{\abs}{\lvert}{\rvert}
\DeclarePairedDelimiterX{\dotp}[2]{\langle}{\rangle}{#1, #2}
\begin{document}
\maketitle

\section{Inner Products, Norms and the Cauchy--Schwarz Inequality}
\label{sec:inner}
Throughout this section, let $V$ be a vector space over the field $\mathbb{R}$.
An \emph{inner product} on $V$ is a map $\langle \cdot, \cdot \rangle \colon V \times V \to \mathbb{R}$ that is linear in each argument, symmetric, and positive definite, meaning that $\langle x, x \rangle \ge 0$ for all $x \in V$, with equality if and only if $x = 0$.

Every inner product induces a \emph{norm} and, through it, a \emph{metric}:
\[
	\norm{x} \coloneqq \sqrt{\dotp{x}{x}}, \qquad d(x, y) \coloneqq \norm{x - y}.
\]
For a nonzero vector $x \neq 0$, we call $\hat{x} \coloneqq x / \norm{x}$ the associated \emph{unit vector}, which satisfies $\norm{\hat{x}} = 1$.

The estimate relating the inner product to the norm is the \emph{Cauchy--Schwarz inequality}
\[
	\abs{\dotp{x}{y}} \le \norm{x} \, \norm{y},
\]
which is satisfied with equality if and only if $x$ and $y$ are linearly dependent.
It guarantees that the angle $\theta$ between two nonzero vectors is well defined through
\[
	\cos\theta = {\dotp{x}{y} \over \norm{x} \, \norm{y}},
\]
and it implies the \emph{triangle inequality} $\norm{x + y} \le \norm{x} + \norm{y}$.
A norm that comes from an inner product is special: it obeys the \emph{parallelogram law}
\[
	\norm{x + y}^2 + \norm{x - y}^2 = 2 \, \norm{x}^2 + 2 \, \norm{y}^2 .
\]

Given an orthonormal system (i.e., a system of vectors that are pairwise normal and of unit length) $e_1, \dots, e_n$, the \emph{orthogonal projection} of $x$ onto its span is
\[
	\textup{proj}_U(x) = \sum_{i=1}^{n} \dotp{x}{e_i} \, e_i ,
\]
where $( x - \textup{proj}_U(x) )$ is orthogonal to every $e_i$.
Comparing the two contributions yields \emph{Bessel's inequality}
\[
	\sum_{i=1}^{n}  \abs{\dotp{x}{e_i}}^2 \le \norm{x}^2 .
\]

\section{The Exponential Function as a Power Series}
\label{sec:exponential}
The real \emph{exponential function} can be defined by the power series
\begin{equation}
\label{eq:exp}
\exp(x) \coloneqq \sum_{n=0}^{\infty} \frac{x^n}{n!},
\end{equation}
which converges for every $x \in \mathbb{R}$. It satisfies $\exp(0) = 1$ and the functional equation $\exp(x + y) = \exp(x)\, \exp(y)$, so in particular $\exp(x)\,\exp(-x) = 1$.

There is also a connection to a classical sequence limit
\begin{equation*}
	\exp(x) = \lim_{n \to \infty} \!\left(1 + \frac{x}{n}\right)^n\!, \qquad\text{giving}\qquad e \coloneqq \exp(1) = \lim_{n \to \infty} \!\Big(1 + \frac{1}{n}\Big)^n.
\end{equation*}

Additionally, the exponential function is its own derivative, 
\[
	\frac{\textup{d}}{\textup{d}x} \exp(x) = \exp(x),
\]
which is very sensible given the original form in  \zcref{eq:exp}, and, together with the initial value $\exp(0) = 1$, characterises it uniquely among all differentiable functions.
\end{document}
